\subsection{La personnalisation de l'apprentissage}
\label{subsec:personnalisation_apprentissage}

L'apprentissage adaptatif désigne l'ajustement du contenu, du rythme et de la méthode pédagogique aux caractéristiques de chaque apprenant. Les systèmes tutoriels classiques implémentaient cette adaptation par des règles explicites : si l'apprenant échoue à une question, présenter un exercice plus simple ; s'il réussit, augmenter la difficulté. Les LLM permettent une adaptation d'une autre nature. Le modèle ajuste non seulement la difficulté, mais aussi la formulation, le style explicatif et le niveau de détail en fonction du contexte conversationnel. Cette adaptation sémantique et stylistique en temps réel constitue une rupture avec les systèmes à règles fixes.

La personnalisation peut également porter sur l'apparence de l'agent. Une stratégie documentée consiste à exploiter la familiarité de l'apprenant avec certaines figures pour activer des schémas relationnels préexistants. L'utilisation de représentations visuelles de personnalités admirées améliore la motivation et la crédibilité perçue de l'agent \citep{pataranutaporn2022}. Le mécanisme repose sur un transfert : la confiance et l'intérêt accordés à la figure connue se reportent sur l'agent qui en adopte l'apparence. Cette stratégie s'inscrit dans la lignée des travaux sur les relations parasociales, où l'exposition médiatique répétée à une personnalité crée un sentiment de familiarité unilatéral.

L'effet est plus marqué lorsque la connexion est personnelle plutôt que parasociale. Une étude comparant un chatbot textuel, une mascotte 3D et un deepfake de l'enseignant réel des élèves montre que la familiarité préexistante avec l'enseignant augmente la confiance et la compétence perçue de l'agent \citep{tan2025}. L'apprenant ne construit pas la crédibilité de l'agent ex nihilo : il transfère celle qu'il accorde déjà à la personne représentée. Ce résultat suggère que l'apparence humaine de l'agent fonctionne comme un signal de disponibilité relationnelle, activant des attentes d'interactivité et de réactivité \citep{linyu2024}.

Cette amélioration de la perception subjective ne se traduit pas systématiquement en gains d'apprentissage objectifs. La distinction entre variables perceptuelles et variables d'apprentissage traverse l'ensemble de la littérature sur les agents pédagogiques (cf.~\ref{subsec:incarnation_agence_sociale}). Les agents génératifs ne font pas exception. Lorsque des études rapportent une équivalence entre agents et instructeurs humains, cette équivalence porte sur des mesures immédiates et des contextes contrôlés \citep{leiker2023, lim2024}. Elle n'indique pas que l'agent reproduit le processus d'enseignement humain ; elle signifie que, pour les indicateurs retenus, les deux conditions produisent des résultats statistiquement indistinguables.

L'apparence humaine de l'agent influence la perception plus que l'apprentissage. Elle agit comme un indice qui signale un potentiel d'interactivité élevé \citep{linyu2024}. L'apprenant anticipe un échange responsif et s'engage en conséquence. Cette dissociation entre perception et apprentissage pose un problème de design : optimiser l'engagement perçu peut se faire au détriment du traitement cognitif profond. L'attrait de l'interface risque de détourner des ressources attentionnelles du contenu vers l'agent, un phénomène déjà documenté avec les instructeurs humains (cf.~\ref{subsec:incarnation_agence_sociale}). Le concepteur doit arbitrer entre deux objectifs potentiellement contradictoires.

La personnalisation par l'apparence soulève également des questions éthiques \citep{danry2022}. L'usage de deepfakes à des fins pédagogiques détourne une technologie associée à la désinformation. Le consentement de la personne représentée, les effets d'une exposition prolongée à des représentations synthétiques et la confusion potentielle entre l'agent et la personne qu'il imite constituent des enjeux que la littérature commence à explorer \citep{berson2024}. Ces questions deviennent particulièrement saillantes lorsque l'agent incarne une figure historique décédée, dont le consentement ne peut être obtenu.

Ces enjeux ne sont pas hypothétiques. Des institutions culturelles majeures déploient déjà des agents conversationnels incarnant des figures disparues. Au Dalí Museum de St. Petersburg, l'installation \textit{Ask Dalí}\footnote{\url{https://thedali.org/exhibit/ask-dali/}} permet aux visiteurs de converser avec une reconstruction de Salvador Dalí via un téléphone-homard, objet iconique de son œuvre surréaliste (voir Figure~\ref{fig:ask_dali}). Le système, alimenté par GPT-4 et une synthèse vocale entraînée sur des archives audio, produit des réponses reflétant la personnalité et l'humour de l'artiste. Au Musée d'Orsay, l'expérience \textit{Hello Vincent}\footnote{\url{https://www.musee-orsay.fr/fr/articles/numerique-bonjour-vincent-275618}} propose une interaction similaire avec Van Gogh, mais s'appuie sur un corpus de près de 900 lettres validé par un historien de l'art (voir Figure~\ref{fig:vangogh_museum}). Ces dispositifs illustrent la transition des prototypes de recherche vers des déploiements grand public. Ils révèlent également une tension : le cadre muséal, orienté vers la médiation culturelle et l'engagement émotionnel, diffère-t-il fondamentalement du cadre scolaire, où l'exactitude factuelle et l'apprentissage conceptuel prévalent ?

\begin{figure}[htbp]
    \centering
    \includegraphics[width=0.7\textwidth]{images/ch2/dali.png}
    \caption{Installation \textit{Ask Dalí} au Dalí Museum de St. Petersburg. Les visiteurs conversent avec une reconstruction IA de Salvador Dalí via un téléphone-homard. La citation murale --- « I do not understand why, when I ask for a grilled lobster in a restaurant, I am never served a cooked telephone » --- illustre l'humour surréaliste que le système cherche à reproduire.}
    \label{fig:ask_dali}
\end{figure}

\begin{figure}[htbp]
    \centering
    \includegraphics[width=0.7\textwidth]{images/ch2/vangogh.png}
    \caption{Représentation numérique de Vincent van Gogh dans l'expérience \textit{Hello Vincent} au Musée d'Orsay. Le visiteur converse avec l'artiste pendant qu'il peint son célèbre \textit{Champ de blé aux corbeaux}, l'IA s'appuyant sur un corpus de près de 900 lettres.}
    \label{fig:vangogh_museum}
\end{figure}

Plusieurs études explorent cette personnification de figures historiques. Des participants adultes conversant avec des reconstructions de Leonardo Da Vinci ou Murasaki Shikibu rapportent une motivation et une efficacité d'apprentissage perçue supérieures à la lecture seule \citep{pataranutaporn2023} (voir Figure~\ref{fig:living_memories}). Des collégiens dialoguant avec Napoléon ou Alexandre le Grand montrent une amélioration de l'empathie historique et des capacités d'auto-réflexion \citep{kim2025}. Ces résultats convergent sur un point : la personnification historique génère un engagement affectif mesurable.

\begin{figure}[htbp]
    \centering
    \includegraphics[width=0.8\textwidth]{images/ch2/living-memories-talking.png}
    \caption{Interface du projet \textit{Living Memories} \citep{pataranutaporn2023}. L'utilisateur converse avec une reconstruction de Leonardo Da Vinci, qui se présente et répond aux questions. Le système combine GPT-3 et un extracteur sémantique pour générer des réponses cohérentes avec le personnage historique.}
    \label{fig:living_memories}
\end{figure}

Cette convergence appelle cependant une lecture critique. Les variables mesurées --- motivation, empathie, engagement --- relèvent du registre perceptuel et affectif. L'apprentissage factuel, la compréhension conceptuelle et la rétention à long terme restent largement inexplorés. L'effet optimal survient lorsque l'interaction complète les sources textuelles plutôt que de s'y substituer \citep{pataranutaporn2023}. Cette observation suggère que la valeur pédagogique réside dans la complémentarité, non dans la substitution. L'agent génératif ne remplace pas l'enseignement : il crée un espace d'engagement dont les effets sur l'apprentissage profond restent à démontrer.

En synthèse, la personnalisation offerte par les agents génératifs opère sur deux registres distincts. L'adaptation sémantique ajuste le contenu aux besoins de l'apprenant et peut soutenir l'apprentissage. L'adaptation par l'apparence exploite des liens relationnels préexistants pour favoriser l'engagement, mais son effet sur l'apprentissage objectif reste à démontrer. La section suivante examine un troisième enjeu : la fiabilité du contenu que ces agents produisent.

